\begin{frame}[shrink]
	\frametitle{Como uma rede neural é treinada}
	\small
	\begin{itemize}
		\item \textbf{1. Forward}: calcula a saída $y$ a partir da entrada e dos pesos atuais.
		\item \textbf{2. Perda (loss)}: compara $y$ com o alvo desejado, num único número que cresce quanto maior o erro.
		\item \textbf{3. Backward}: a regra da cadeia calcula $\partial L/\partial w$ --- o quanto a perda mudaria se cada peso mudasse um pouco.
		\item \textbf{4. Atualização}: cada peso anda um pequeno passo no sentido \textbf{contrário} ao gradiente.
	\end{itemize}
	\begin{equation*}
		w \leftarrow w - \eta \cdot \frac{\partial L}{\partial w}
	\end{equation*}
	\par Exemplo numérico real (neurônio único, $x=2{,}0$, alvo $=0{,}9$, $w_0=-1{,}0$, taxa de aprendizado $\eta=3{,}0$):
	\begin{equation*}
		z = w_0 x = -2{,}0 \ \Rightarrow\ y = \sigma(z) \approx 0{,}1192 \ \Rightarrow\ L = \tfrac12(y-\text{alvo})^2 \approx 0{,}3048
	\end{equation*}
	\begin{equation*}
		\frac{\partial L}{\partial w} \approx -0{,}1640
		\quad\Longrightarrow\quad
		w_1 = -1{,}0 - 3{,}0 \cdot(-0{,}1640) \approx -0{,}508
	\end{equation*}
	\par Repetindo esses quatro passos, iteração após iteração, $y$ converge para o alvo. É exatamente esse ciclo --- forward, perda, backward, atualização --- que o resto desta palestra questiona: o que acontece quando os pesos só podem valer $\{-1,0,1\}$ (BitNets), ou quando a ativação vira um pulso discreto no tempo (redes de pulso)?
	% [SOFTWARE] Backprop -> Forward e backward classicos | passo a passo dos
	% 9 estagios da iteracao 1 reproduz exatamente estes numeros
	\abrirNoSoftware{Backprop \textrightarrow{} Forward e backward clássicos}{backprop.classic}
\end{frame}
